\documentclass[11pt]{article}
\usepackage[utf8]{inputenc}
\usepackage[spanish]{babel}
\usepackage{amsmath,amssymb}
\usepackage{booktabs}
\usepackage[margin=2.5cm]{geometry}
\setlength{\parindent}{0pt}
\setlength{\parskip}{6pt}

\begin{document}
\begin{center}
	{\small SMAII, Redes Complejas Multicapa}\\[6pt]
	{\Large\bfseries Actividad 2}\\[3pt]
	{\small Octubre de 2026}
\end{center}

\textbf{Integrantes del equipo}
\begin{itemize}
    \item Jorge Miyoshi Mikami Alonso
    \item Luis Manuel Reyes Colín 
\end{itemize}

De la figura se obtienen las siguientes aristas dirigidas ($A \to B$ con peso $=$ pases):

\begin{center}
\begin{tabular}{cll}
\toprule
Nodo & Aristas (peso) & Exgrado \\
\midrule
1 & 2 (5), 3 (2), 4 (1) & 3 \\
2 & 3 (4), 5 (2)        & 2 \\
3 & 1 (3), 6 (2), 9 (3) & 3 \\
4 & 2 (5), 5 (2), 7 (1) & 3 \\
5 & 4 (4), 8 (6)        & 2 \\
6 & 5 (2), 9 (2)        & 2 \\
7 & 6 (9), 8 (2)        & 2 \\
8 & 7 (3)               & 1 \\
9 & 8 (3)               & 1 \\
\bottomrule
\end{tabular}
\end{center}

\textbf{a)}

En la Katz normal, un nodo es importante si le \emph{llegan} pases de nodos importantes (ingrado). En la Katz de exgrado es al revés, o sea, un nodo es importante si \emph{da} pases a nodos importantes.

La centralidad de Katz es
\[
C = (I - \alpha A)^{-1}\,\beta\mathbf{1}, \qquad \beta = 1 .
\]
La inversa se puede escribir como una suma ($\tfrac{1}{1-x} = 1 + x + x^2 + \cdots$):
\[
(I - \alpha A)^{-1} = I + \alpha A + \alpha^2 A^2 + \alpha^3 A^3 + \cdots
\]
Entonces
\[
C = \mathbf{1} + \alpha\,A\mathbf{1} + \alpha^2 A^2\mathbf{1} + \alpha^3 A^3\mathbf{1} + \cdots
\]

Como $A_{ij}=1$ si $i$ le pasa a $j$ (son adyacentes), interpretamos: 
\begin{itemize}
  \item $(A\mathbf{1})_i$ = número de caminos de longitud 1 que \emph{salen} de $i$ = exgrado de $i$.
  \item $(A^2\mathbf{1})_i = \sum_{j:\, i\to j} (A\mathbf{1})_j$ = caminos de longitud 2 que salen de $i$
        (se suman los exgrados de los nodos a los que $i$ pasa).
  \item $(A^3\mathbf{1})_i = \sum_{j:\, i\to j} (A^2\mathbf{1})_j$ = caminos de longitud 3 que salen de $i$.
\end{itemize}

Por ejemplo, para el nodo 1 (pasa a 2, 3 y 4):
\[
(A^2\mathbf{1})_1 = 2 + 3 + 3 = 8, \qquad (A^3\mathbf{1})_1 = 5 + 6 + 6 = 17 .
\]

Haciendo lo mismo para todos los nodos, considerando hasta el tercer nivel:

\begin{center}
\begin{tabular}{ccccc}
\toprule
Nodo & Nivel 1: $A\mathbf{1}$ & Nivel 2: $A^2\mathbf{1}$ & Nivel 3: $A^3\mathbf{1}$ & $C$ con $\alpha=0.1$ \\
\midrule
1 & 3 & 8 & 17 & 1.397 \\
2 & 2 & 5 & 10 & 1.260 \\
3 & 3 & 6 & 12 & 1.372 \\
4 & 3 & 6 & 12 & 1.372 \\
5 & 2 & 4 & 8  & 1.248 \\
6 & 2 & 3 & 5  & 1.235 \\
7 & 2 & 3 & 5  & 1.235 \\
8 & 1 & 2 & 3  & 1.123 \\
9 & 1 & 1 & 2  & 1.112 \\
\bottomrule
\end{tabular}
\end{center}

(Se usa $\alpha = 0.1$ solo como ejemplo)

\textbf{Resultado:}
\begin{enumerate}
  \item \textbf{Nodo 1}: tiene el exgrado máximo (3), empatado con 3 y 4, pero sus pases llegan a
        nodos que también pasan mucho (2, 3 y 4), así que gana en el nivel 2 (8) y en el nivel 3 (17).
  \item[2--3.] \textbf{Nodos 3 y 4}: también tienen exgrado 3 y quedan empatados
        en los niveles 2 (6) y 3 (12). 
\end{enumerate}
\textbf{Nota:} Para deshacer el empate, calculamos un nivel más:
\[
(A^4\mathbf{1})_3 = 24 > 23 = (A^4\mathbf{1})_4.
\]

Así que el nodo 3 queda ligeramente por encima del nodo 4.
Por lo tanto, el orden es $1 > 3 \gtrsim 4$. Visualmente, los tres son los nodos de los que salen
más flechas y cuyas flechas llegan a nodos que a su vez tienen muchas flechas de salida.

\textbf{b)}

Un nodo con alta centralidad de Katz de exgrado es un jugador que da muchos pases y además se los da a 
jugadores que también pasan el balón. Entonces, el balón puede llegar a muchos jugadores a partir de él 
en pocos pases. En fútbol sería un jugador que genera o organiza el juego, como un mediocampista.

\textbf{c)}

La centralidad de ingrado es local porque solo cuenta las aristas que llegan directamente al nodo. Todas cuentan igual, sin importar de qué nodo vienen.

La centralidad de Katz es global, porque también toma en cuenta los caminos indirectos que llegan al nodo. Mientras más largo sea el camino, menos pesa por $\alpha^k$. Además, recibir conexiones de nodos importantes hace que el nodo tenga más importancia. El término $\beta$ hace que todos los nodos tengan una importancia mínima, aunque no reciban ninguna conexión. Por eso, el ingrado se puede ver como quedarse solo con el primer término de Katz.

\end{document}
